%═══════════════════════════════════════════════════════════════════════════════
\chapter{Teoría de la Información y Mecánica Estadística}
\label{ap:info_mecanica}
%═══════════════════════════════════════════════════════════════════════════════

Este apéndice es el argumento central del libro expuesto en una sola pieza.
La tesis —que el aprendizaje automático moderno es física estadística en
dimensiones altas— descansa sobre un conjunto de identidades matemáticas
precisas, no sobre una analogía. El objetivo aquí es presentar esas identidades
de forma autocontenida, de modo que el lector pueda leer este apéndice
independientemente del resto del texto y entender por qué la tesis es verdadera.

Cada sección indica los capítulos del libro donde el resultado aparece en uso.

%───────────────────────────────────────────────────────────────────────────────
\section{La identidad entre entropía física y entropía de la información}
\label{sec:ap_entropia}
%───────────────────────────────────────────────────────────────────────────────

\textit{Usado en: Capítulo~\ref{ch:cap7} (entropía como medida de incertidumbre),
Capítulo~\ref{ch:cap11} (H-teorema), Capítulo~\ref{ch:cap12} (ELBO como
energía libre).}

\medskip
Hay dos entropías en la ciencia moderna. La primera es la de Boltzmann (1877):
\begin{equation}
  S = k_B\log\Omega,
  \label{eq:ap_boltzmann}
\end{equation}
donde $\Omega$ es el número de microestados compatibles con un macroestado dado.
La segunda es la de Shannon (1948): para una distribución discreta
$\{p_i\}_{i=1}^n$,
\begin{equation}
  H = -\sum_{i=1}^n p_i\log p_i.
  \label{eq:ap_shannon}
\end{equation}
Son la misma función.

Para verlo, sea un sistema con $\Omega$ microestados equiprobables. Cada uno
tiene probabilidad $p_i = 1/\Omega$. La entropía de Shannon es:
\[
  H = -\sum_{i=1}^\Omega \frac{1}{\Omega}\log\frac{1}{\Omega} = \log\Omega.
\]
Salvo el factor $k_B$ (que es una elección de unidades, no un contenido físico),
$H = S/k_B$. La entropía de Boltzmann es el caso especial de la entropía de
Shannon cuando todos los microestados son equiprobables. La generalización de
Shannon a distribuciones arbitrarias permite hablar de entropía para sistemas
fuera del equilibrio, con distribuciones no uniformes.

Para densidades continuas $p(\bm{x})$ sobre $\mathbb{R}^d$, la
\textbf{entropía diferencial} es:
\begin{equation}
  H[p] = -\int p(\bm{x})\log p(\bm{x})\,d\bm{x} = -\mathbb{E}_p[\log p].
  \label{eq:ap_entropia_cont}
\end{equation}

\begin{proposicion}{Entropía máxima bajo restricciones de momento}
\label{prop:ap_maxent}
La distribución de máxima entropía sobre $\mathbb{R}^d$ sujeta a
$\mathbb{E}[\bm{x}] = \bm{\mu}$ y $\mathbb{E}[\bm{x}\bm{x}^\top] = \bm{\Sigma}
+ \bm{\mu}\bm{\mu}^\top$ (media y covarianza fijas) es la gaussiana
$\mathcal{N}(\bm{\mu},\bm{\Sigma})$.
\end{proposicion}

\begin{proof}
Por multiplicadores de Lagrange: maximizar $H[p] = -\int p\log p$ sujeto a
$\int p = 1$, $\int p\,x_i = \mu_i$ y $\int p\,x_ix_j = \Sigma_{ij}+\mu_i\mu_j$.
La condición de optimalidad $\delta H/\delta p = 0$ da
$\log p(\bm{x}) = -1 + \lambda_0 + \bm{\lambda}_1^\top\bm{x} +
\bm{x}^\top\bm{\Lambda}_2\bm{x}$, que es cuadrático en $\bm{x}$: precisamente
una gaussiana. Los multiplicadores $\lambda_0, \bm{\lambda}_1, \bm{\Lambda}_2$
se determinan imponiendo las restricciones de momento. $\square$
\end{proof}

Esta proposición explica por qué la gaussiana aparece en tantos contextos: no
es un supuesto arbitrario sino la distribución de máxima incertidumbre dado
el conocimiento de media y varianza. Es la distribución más honesta posible
cuando solo se conocen esos momentos.

%───────────────────────────────────────────────────────────────────────────────
\section{El principio de máxima entropía de Jaynes}
\label{sec:ap_jaynes}
%───────────────────────────────────────────────────────────────────────────────

\textit{Usado en: Capítulo~\ref{ch:cap7} (distribución de Gibbs como máxima
entropía, temperatura como multiplicador de Lagrange), Capítulo~\ref{ch:cap12}
(CFG scale como temperatura inversa).}

\medskip
Jaynes (1957) reformuló la mecánica estadística como un problema de inferencia:
dada la información disponible sobre un sistema, la distribución de probabilidad
que asignarle es aquella que maximiza la entropía sujeta a las restricciones
que impone la información conocida. El resultado no es un modelo del sistema
físico sino el estado de conocimiento más honesto compatible con los datos.

\begin{teorema}{Distribución de Gibbs-Boltzmann como máxima entropía}
\label{thm:ap_jaynes}
La distribución de máxima entropía sobre $\mathbb{R}^d$ sujeta a
$\int p(\bm{x})\,d\bm{x} = 1$ y $\mathbb{E}_p[U(\bm{x})] = \langle U\rangle$
(valor medio de la energía $U$ fijo) es la \textbf{distribución de
Gibbs-Boltzmann}:
\begin{equation}
  p^*(\bm{x}) = \frac{1}{Z(\beta)}\,e^{-\beta U(\bm{x})},
  \qquad Z(\beta) = \int e^{-\beta U(\bm{x})}\,d\bm{x},
  \label{eq:ap_gibbs}
\end{equation}
donde $\beta$ es el multiplicador de Lagrange asociado a la restricción de
energía media, y $Z(\beta)$ es la \textbf{función de partición}.
\end{teorema}

\begin{proof}
Funcional de Lagrange con multiplicadores $\alpha$ (normalización) y $\beta$
(energía media):
\[
  \frac{\delta}{\delta p}\left[-\int p\log p - \alpha\!\int p - \beta\!\int pU\right]
  = -\log p - 1 - \alpha - \beta U = 0.
\]
Luego $p^*(\bm{x}) = e^{-1-\alpha}e^{-\beta U(\bm{x})} =: Z^{-1}e^{-\beta U(\bm{x})}$,
donde $Z = e^{1+\alpha}$ se determina por la normalización. $\square$
\end{proof}

El multiplicador de Lagrange $\beta$ tiene interpretación física directa. De la
termodinámica: $\beta = 1/(k_BT)$, donde $T$ es la temperatura. Cuanto mayor es
$T$ (menor $\beta$), más difusa es la distribución; cuanto menor es $T$ (mayor
$\beta$), más concentrada en los mínimos de $U$.

\medskip\noindent\textbf{La función de partición como objeto central.}
Todo se deriva de $Z(\beta)$:
\begin{align}
  \text{Energía libre de Helmholtz:}&\quad
    F(\beta) = -\frac{1}{\beta}\log Z(\beta). \label{eq:ap_helmholtz}\\
  \text{Energía media:}&\quad
    \langle U\rangle = -\frac{\partial\log Z}{\partial\beta}. \label{eq:ap_U}\\
  \text{Entropía:}&\quad
    S = k_B(\log Z + \beta\langle U\rangle) = k_B\beta^2\frac{\partial F}{\partial\beta}.
    \label{eq:ap_S}
\end{align}
La relación fundamental $F = \langle U\rangle - TS$ (en unidades donde $k_B=1$)
conecta la energía libre con la energía y la entropía. En el
Capítulo~\ref{ch:cap12}, el ELBO de DDPM es precisamente $-F$: maximizar el
ELBO equivale a minimizar la energía libre variacional.

%───────────────────────────────────────────────────────────────────────────────
\section{Energía libre variacional y el ELBO}
\label{sec:ap_elbo_helmholtz}
%───────────────────────────────────────────────────────────────────────────────

\textit{Usado en: Capítulo~\ref{ch:cap7} (VAE como minimización de energía libre),
Capítulo~\ref{ch:cap12} (DDPM como energía libre variacional).}

\medskip
La función de partición $Z = \int e^{-\beta U(\bm{x})}\,d\bm{x}$ no es computable
directamente cuando $U$ es compleja (alta dimensión, múltiples mínimos). La
aproximación variacional de la energía libre introduce una distribución
$q(\bm{x})$ tratable y acota $F$ por abajo.

\begin{proposicion}{Energía libre variacional}
\label{prop:ap_flv}
Para cualquier distribución $q(\bm{x})$ con soporte en el soporte de $p^*$:
\begin{equation}
  F = -\frac{1}{\beta}\log Z
  \leq \mathbb{E}_q[U(\bm{x})] - \frac{1}{\beta}H[q]
  =: F_q,
  \label{eq:ap_flv}
\end{equation}
con igualdad si y solo si $q = p^*$.
\end{proposicion}

\begin{proof}
$F_q - F = \mathbb{E}_q[U] + \frac{1}{\beta}\mathbb{E}_q[\log q] + \frac{1}{\beta}
\log Z = \frac{1}{\beta}\mathbb{E}_q[\log q + \beta U + \log Z]
= \frac{1}{\beta}\mathbb{E}_q[\log(q/p^*)] = \frac{1}{\beta}D_{\mathrm{KL}}(q\|p^*)
\geq 0$.
La igualdad requiere $q = p^*$. $\square$
\end{proof}

La energía libre variacional $F_q = \langle U\rangle_q - T H[q]$ tiene dos
términos que compiten: el primer término minimiza la energía (favorece regiones
de bajo $U$) y el segundo maximiza la entropía (favorece distribuciones difusas).
El balance entre ellos depende de la temperatura $T = 1/\beta$.

\medskip\noindent\textbf{Traducción al ELBO del Capítulo~\ref{ch:cap7}.}
Identificando $U(\bm{z}) = -\log p_\theta(\bm{x}|\bm{z}) - \log p(\bm{z})$
(la energía negativa del modelo generativo) y $q = q_\phi(\bm{z}|\bm{x})$
(el encoder del VAE), con $\beta = 1$:
\begin{align*}
  F_q &= \mathbb{E}_{q_\phi}[-\log p_\theta(\bm{x}|\bm{z}) - \log p(\bm{z})]
         + \mathbb{E}_{q_\phi}[\log q_\phi(\bm{z}|\bm{x})] \\
      &= -\mathbb{E}_{q_\phi}[\log p_\theta(\bm{x}|\bm{z})]
         + D_{\mathrm{KL}}(q_\phi(\bm{z}|\bm{x})\|p(\bm{z})).
\end{align*}
Minimizar $F_q$ es equivalente a maximizar el ELBO:
$\mathrm{ELBO} = -F_q = \mathbb{E}_{q_\phi}[\log p_\theta(\bm{x}|\bm{z})]
- D_{\mathrm{KL}}(q_\phi\|p)$.
El primer término es la energía de reconstrucción; el segundo es la entropía
relativa que regulariza el espacio latente. El VAE minimiza exactamente la
energía libre de Helmholtz.

%───────────────────────────────────────────────────────────────────────────────
\section{La ecuación de Langevin y la distribución de Gibbs del SGD}
\label{sec:ap_langevin_gibbs}
%───────────────────────────────────────────────────────────────────────────────

\textit{Usado en: Capítulo~\ref{ch:cap5} (SDE del SGD, distribución estacionaria),
Capítulo~\ref{ch:cap11} (Proposición sobre la distribución de Gibbs,
decrecimiento de la KL).}

\medskip
La ecuación de Langevin es la SDE de un sistema físico bajo la acción de una
fuerza derivada de un potencial $U(\bm{x})$ y ruido térmico:
\begin{equation}
  d\bm{x}_t = -\nabla U(\bm{x}_t)\,dt + \sqrt{2k_BT}\,d\bm{W}_t.
  \label{eq:ap_langevin}
\end{equation}
El coeficiente del ruido es $\sqrt{2k_BT}$ por la relación de fluctuación-disipación:
la misma temperatura que determina la amplitud del ruido determina la distribución
estacionaria. La distribución de Gibbs~\eqref{eq:ap_gibbs} con $\beta = 1/(k_BT)$
y potencial $U = \mathcal{L}$ es la distribución estacionaria de esta SDE.

La SDE del SGD del Capítulo~\ref{ch:cap5} es la ecuación de Langevin con:
\[
  U(\bm{\theta}) = \mathcal{L}(\bm{\theta}), \qquad k_BT = \frac{\eta c}{2},
\]
donde $\mathcal{L}$ es la función de pérdida, $\eta$ la tasa de aprendizaje y
$c$ la varianza del gradiente estocástico. La distribución estacionaria del SGD
es la distribución de Gibbs $p^*\propto e^{-2\mathcal{L}/(\eta c)}$: el SGD
explora el paisaje de pérdida como un sistema físico a temperatura efectiva
$T_{\mathrm{eff}} = \eta c/2$.

\begin{observacion}{La tabla de correspondencias Langevin--SGD}
\begin{center}
\begin{tabular}{ll}
\toprule
\textbf{Física estadística} & \textbf{Aprendizaje automático} \\
\midrule
Potencial $U(\bm{x})$ & Función de pérdida $\mathcal{L}(\bm{\theta})$ \\
Temperatura $k_BT$ & Temperatura efectiva $\eta c/2$ \\
Fuerza $-\nabla U$ & Gradiente negativo $-\nabla\mathcal{L}$ \\
Distribución de Gibbs $e^{-U/(k_BT)}/Z$ & Distribución estacionaria del SGD \\
Energía libre $F = \langle U\rangle - TS$ & ELBO del modelo variacional \\
Recocido simulado: $T\to 0$ & Learning rate decay: $\eta\to 0$ \\
Mínimo de energía global & Parámetros óptimos $\bm{\theta}^*$ \\
\bottomrule
\end{tabular}
\end{center}
\end{observacion}

%───────────────────────────────────────────────────────────────────────────────
\section{El H-teorema de Boltzmann y el decrecimiento de la KL}
\label{sec:ap_hteorema}
%───────────────────────────────────────────────────────────────────────────────

\textit{Usado en: Capítulo~\ref{ch:cap11} (decrecimiento de $D_{\mathrm{KL}}$
hacia el equilibrio, tasa de producción de entropía).}

\medskip
El H-teorema de Boltzmann (1872) afirma que la función $H(t) = \int p_t\log p_t$
decrece monótonamente hacia su mínimo en el equilibrio. En el lenguaje moderno,
es el enunciado de que la divergencia KL respecto a la distribución de equilibrio
decrece a lo largo de la dinámica.

\begin{teorema}{H-teorema cuantitativo}
\label{thm:ap_h}
Para la SDE de Langevin~\eqref{eq:ap_langevin} con distribución estacionaria
$p^*\propto e^{-\beta U}$ y distribución marginal $p_t$ en el tiempo $t$:
\begin{equation}
  \frac{d}{dt}D_{\mathrm{KL}}(p_t\|p^*)
  = -k_BT\,\mathcal{I}(p_t\|p^*),
  \label{eq:ap_h_cuantitativo}
\end{equation}
donde $\mathcal{I}(p_t\|p^*) = \mathbb{E}_{p_t}[\|\nabla\log p_t -
\nabla\log p^*\|^2]$ es la \textbf{información de Fisher relativa}
(o disipación de entropía).
\end{teorema}

\begin{proof}
Sea $\phi_t = \log(p_t/p^*)$. Por la ecuación de Fokker-Planck
(Capítulo~\ref{ch:cap11}, ecuación~\eqref{eq:fp}) con $\bm{\mu} = -\nabla U$
y $\bm{D} = k_BT\bm{I}$:
\begin{align*}
  \frac{d}{dt}D_{\mathrm{KL}}(p_t\|p^*)
  &= \int\partial_t p_t\,\phi_t\,d\bm{x} \\
  &= \int\phi_t\left[\nabla\cdot(\nabla U\,p_t) + k_BT\Delta p_t\right]\,d\bm{x}.
\end{align*}
Integrando por partes dos veces y usando que $p^*\propto e^{-\beta U}$ (luego
$\nabla\log p^* = -\beta\nabla U = -\nabla U/(k_BT)$):
\begin{align*}
  \frac{d}{dt}D_{\mathrm{KL}}
  &= -\int p_t\,\nabla\phi_t\cdot(\nabla U + k_BT\nabla\phi_t)\,d\bm{x}\\
  &= -k_BT\int p_t\|\nabla\phi_t\|^2\,d\bm{x}\\
  &= -k_BT\,\mathbb{E}_{p_t}[\|\nabla\log p_t - \nabla\log p^*\|^2].
\end{align*}
Donde usamos $\nabla U = -k_BT\nabla\log p^*$ y que el término cruzado
$-\int p_t\nabla\phi_t\cdot\nabla U\,d\bm{x} = k_BT\int p_t\nabla\phi_t\cdot
\nabla\log p^*\,d\bm{x}$ se cancela al completar el cuadrado. $\square$
\end{proof}

La información de Fisher relativa $\mathcal{I}(p_t\|p^*)$ mide cuánto difieren
los campos vectoriales $\nabla\log p_t$ y $\nabla\log p^*$: en el equilibrio
$p_t = p^*$, ambos coinciden y $\mathcal{I} = 0$, consistente con $d(\mathrm{KL})/dt
= 0$. Fuera del equilibrio, la disipación es positiva y la KL decrece. Este es
el H-teorema de Boltzmann en forma exacta.

\medskip\noindent\textbf{Conexión con el score del Capítulo~\ref{ch:cap11}.}
El campo $\nabla\log p_t$ es exactamente el score $\bm{s}_t$ de los modelos de
difusión. La ecuación~\eqref{eq:ap_h_cuantitativo} dice que cuanto más se aleja
el score actual del score de equilibrio, más rápido decrece la KL. El reverse
SDE de Anderson mueve el sistema en la dirección de mayor disipación.

%───────────────────────────────────────────────────────────────────────────────
\section{Entropía mutua, información y el cuello de botella de la información}
\label{sec:ap_info_mutua}
%───────────────────────────────────────────────────────────────────────────────

\textit{Usado en: Capítulo~\ref{ch:cap7} (representaciones internas de las redes
y el principio del cuello de botella de la información).}

\medskip
La \textbf{entropía mutua} entre dos variables $X$ e $Y$ mide cuánta información
comparten:
\begin{equation}
  I(X;Y) = D_{\mathrm{KL}}(p(X,Y)\|p(X)p(Y))
  = H(X) - H(X|Y) = H(Y) - H(Y|X),
  \label{eq:ap_mutua}
\end{equation}
donde $H(X|Y) = -\mathbb{E}_{p(X,Y)}[\log p(X|Y)]$ es la entropía condicional.

La entropía mutua es simétrica: $I(X;Y) = I(Y;X)$. Es cero si y solo si $X$
e $Y$ son independientes. Su interpretación: $I(X;Y)$ es la reducción promedio
de incertidumbre sobre $X$ al conocer $Y$.

\subsection*{La cadena de Markov y las desigualdades de procesado de datos}

Si $X \to Y \to Z$ forman una cadena de Markov ($Z$ es condicionalmente
independiente de $X$ dado $Y$):
\begin{equation}
  I(X;Z) \leq I(X;Y), \qquad I(X;Z) \leq I(Y;Z).
  \label{eq:ap_dpi}
\end{equation}
Estas son las \textbf{desigualdades de procesado de datos}: aplicar cualquier
función a los datos no puede crear información nueva. Cada capa de una red
neuronal es un procesamiento que solo puede reducir la información sobre la
entrada.

\subsection*{El principio del cuello de botella de la información}

El principio de Tishby et al.\ (2000) afirma que una representación interna
$T$ de una red para la tarea de predecir $Y$ a partir de $X$ óptima maximiza
$I(T;Y)$ sujeta a minimizar $I(T;X)$: comprime la información irrelevante
sobre $X$ (la que no predice $Y$) mientras preserva la relevante.

El problema variacional correspondiente es:
\begin{equation}
  \min_{p(T|X)}\; I(T;X) - \beta\, I(T;Y),
  \label{eq:ap_ib}
\end{equation}
donde $\beta > 0$ controla el balance compresión-preservación. Para $\beta$
grande, la representación preserva casi toda la información relevante; para
$\beta$ pequeño, comprime agresivamente. El parámetro $\beta$ es la temperatura
del Capítulo~\ref{ch:cap5}: el cuello de botella de la información es la
energía libre variacional del espacio de representaciones.

%───────────────────────────────────────────────────────────────────────────────
\section{El dualismo completo: la tabla de correspondencias del libro}
\label{sec:ap_tabla_completa}
%───────────────────────────────────────────────────────────────────────────────

El siguiente cuadro condensa las identidades matemáticas que sostienen la tesis
del libro. No son analogías: son igualdades.

\begin{center}
\begin{tabular}{p{6.2cm}p{6.2cm}}
\toprule
\textbf{Física estadística} & \textbf{Aprendizaje automático} \\
\midrule
Entropía de Boltzmann $S = k_B\log\Omega$ &
Entropía de Shannon $H = -\sum p_i\log p_i$ \\[4pt]

Distribución de Gibbs $p^*\propto e^{-\beta U}$ &
Distribución estacionaria del SGD $p^*\propto e^{-2\mathcal{L}/(\eta c)}$ \\[4pt]

Temperatura $k_BT$ &
Temperatura efectiva del SGD $\eta c/2$ \\[4pt]

Energía libre de Helmholtz $F = \langle U\rangle - TS$ &
ELBO $= -\langle\text{reconstrucción}\rangle + D_\mathrm{KL}$ \\[4pt]

Principio de máx.\ entropía de Jaynes &
Inferencia bayesiana variacional \\[4pt]

H-teorema: $\dot{H}\leq 0$ en la dinámica de Langevin &
Decrecimiento de la KL hacia $p^*$ \\[4pt]

Información de Fisher $\mathcal{I}$ como tasa de disipación &
Gradiente natural = preacondicionamiento por $\bm{I}^{-1}$ \\[4pt]

Función de partición $Z = \int e^{-\beta U}$ &
Evidencia $p(\bm{x}) = \int p(\bm{x}|\bm{z})p(\bm{z})\,d\bm{z}$ \\[4pt]

Score $\nabla\log p$ como fuerza termodinámica $-\nabla U/(k_BT)$ &
Score $\nabla\log p_t$ en el reverse SDE de Anderson \\[4pt]

Recocido simulado: $T(t)\to 0$ &
Learning rate schedule: $\eta(t)\to 0$ \\[4pt]

Temperatura inversa $\beta = 1/(k_BT)$ &
Guidance scale $s$ en CFG \\[4pt]

Puente de Schrödinger: proceso de mínima KL con marginales fijas &
Proceso reverso de Anderson como OT óptimo \\
\bottomrule
\end{tabular}
\end{center}

\medskip
Cada fila de esta tabla es una identidad matemática demostrada en el texto o
en este apéndice, no una metáfora. La diferencia entre la columna izquierda y
la derecha es el dominio de aplicación y la escala dimensional, no la matemática.

%───────────────────────────────────────────────────────────────────────────────
\begin{notasbib}
La entropía de Shannon fue introducida en Shannon (1948)~\cite{shannon1948mathematical}.
La conexión entre entropía de Shannon y entropía de Boltzmann, y el principio de
máxima entropía como fundamento de la mecánica estadística, es el programa de
Jaynes (1957)~\cite{jaynes1957information}. La presentación más completa del
programa de Jaynes está en Jaynes (2003)~\cite{jaynes2003probability}.

La teoría de la información en el sentido moderno, incluyendo la entropía mutua
y las desigualdades de procesado de datos, está en Cover y Thomas
(2006)~\cite{cover2006elements}: el texto de referencia estándar del campo. El
principio del cuello de botella de la información es de Tishby, Pereira y Bialek
(2000)~\cite{tishby2000information}, citado en el Capítulo~\ref{ch:cap7}.

La energía libre variacional y su relación con el ELBO de los modelos de variable
latente están en Kingma y Welling (2013)~\cite{kingma2013auto}, con la conexión
con la termodinámica discutida en detalle en Hinton y Zemel (1994), que
introdujeron el nombre ``energía libre'' en el contexto del aprendizaje
automático.

El H-teorema de Boltzmann en su forma cuantitativa~\eqref{eq:ap_h_cuantitativo},
expresando la tasa de decrecimiento de la KL como la información de Fisher
relativa, es un resultado clásico de la teoría de procesos de difusión. Está
demostrado en Risken (1989)~\cite{risken1989fokker} para la ecuación de
Fokker-Planck, y en el contexto de la teoría de la información en
Villani (2009)~\cite{villani2009optimal}. La conexión con el score de los
modelos de difusión sigue a Song et al.\ (2021)~\cite{song2021score}.

La identidad entre el gradiente natural y el preacondicionamiento por la
información de Fisher, mencionada en la tabla de la Sección~\ref{sec:ap_tabla_completa},
es de Amari (1998)~\cite{amari1998natural}, citado en el Capítulo~\ref{ch:cap7}.
\end{notasbib}
